\subsection{Algebraic points on the actual quotient graph}

\begin{result}{A.4.38}{Arithmetic rank of the conjugate quotient graph}{Proved}
Suppose the actual coefficient $h$ of $E/\Qb$ is algebraic, and
retain the actual scale $c$ in A.4.33. The abelian group
\[
 S_h=\operatorname{graph}(\Phi_h)\cap
                 (E^\dagger\times\bar E)(\Qb),\qquad
 \Phi_h[(u,s)]=[(s-hu)/c],
\]
has torsion-free rank at most two. The two universal-extension
periods alone satisfy the numerical conclusion of the subgroup
theorem used here and give no contradiction to this hypothesis.
\end{result}

\begin{proof}
Put $G=E^\dagger\times\bar E$ and
\[
 V=\{(u,s,v):s-hu=cv\}\subset\operatorname{Lie}G,\qquad
 W=\mathbf C(1,h,0)\subset V.
\]
Here $\dim G=3$, $\dim V=2$, and $W$ is defined over $\Qb$.
Its exponential image is the actual graph. The intersection of
$V$ with the period kernel is exactly
\[
 \{(\omega,\eta_\omega,\bar\omega):\omega\in\Lambda\},
\]
which has rank two: the equation defining $V$ uniquely determines
the third coordinate of a source period. This retains $c$ and does
not rescale it algebraically.

We use \cite[Theorem 4.1, pp.385--386]{EllWaldschmidt1988Subgroup}
with its exact rank convention. For a finitely generated
$Y\subset V$ satisfying $\Gamma=\exp_G(Y)\subset G(\Qb)$, put
$\kappa=\operatorname{rank}(Y\cap\ker\exp_G)$. The theorem supplies
a proper connected algebraic subgroup $G'$ such that, with
\[
 \delta=\dim(G/G'),\quad
 t=\dim\bigl(W/(W\cap\operatorname{Lie}G')\bigr),\quad
 \lambda=\operatorname{rank}\bigl(\Gamma/(\Gamma\cap G')\bigr),
\]
one has $\delta>t$ and, in the present case,
\[
 (\lambda+2\delta)(3-2)\le(\delta-t)(6-\kappa).
\]
Here the theorem's direct additive and multiplicative factors have
dimensions zero, and its remaining factor is all of $G$.
In particular $\lambda$ is the rank \emph{after exponentiation};
the period rank is the separate term $\kappa$.
For $Y$ consisting just of the two period generators,
$\Gamma=0$, $\kappa=2$, and $G'=0$ gives $6\le8$ and $3>1$.
This verifies the asserted period-only scope exactly.

The kernel of $\Phi_h$ is $\exp(\mathbf C(1,h))$:
if $(s-hu)/c\in\bar\Lambda$, subtract the corresponding unique
source period to obtain a vector in this line. Its only algebraic
point is the identity. Otherwise the analytic subgroup theorem
\cite[Hauptsatz and Bemerkung 2, printed p.3]{EllWustholz1984}
would give a positive-dimensional algebraic subgroup with Lie
algebra contained in that line, contrary to A.4.15.
Therefore
\begin{equation}\label{ell:graph-kernel-algebraic}
 S_h\cap(E^\dagger\times\{0\})=\{0\}.
\end{equation}

Now take any finitely generated $\Gamma\subset S_h$, choose
logarithms of its generators in $V$, and adjoin the two period
generators. Their span $Y$ has $\kappa=2$, so the inequality becomes
\begin{equation}\label{ell:graph-rank-inequality}
 \lambda\le2\delta-4t,\qquad\delta>t.
\end{equation}
We check all proper connected subgroups that can satisfy it.
A one-dimensional subgroup is either the vertical
$\mathbf G_a\times\{0\}$ or $\{0\}\times\bar E$.
Indeed an affine subgroup maps trivially to both elliptic factors;
a torus also maps trivially to the remaining additive factor.
An elliptic subgroup with nonzero projection to $E^\dagger$ would
split the nonzero universal extension after an isogeny. This is
impossible because isogeny pullback on $H^1(\mathcal O_E)$ is
injective in characteristic zero, as its trace composed with
pullback is multiplication by its degree.
Neither possible subgroup contains $W$, and each meets $\Gamma$
trivially, by \eqref{ell:graph-kernel-algebraic} and the graph
property. Thus $\delta=2$, $t=1$, $\lambda=\operatorname{rank}\Gamma$,
and \eqref{ell:graph-rank-inequality} gives rank zero.

For a two-dimensional $G'$, the same inequality requires $t=0$.
Let $q:G\to E\times\bar E$ be the quotient by the vertical
$\mathbf G_a$. If $q(G')$ has dimension one, its Lie algebra contains
$\operatorname{Lie}E\times0$, since $W\subset\operatorname{Lie}G'$.
Its projection to $\bar E$ has zero differential and is trivial in
characteristic zero. Hence $q(G')=E\times0$ and
$G'=E^\dagger\times0$. If $q(G')$ has dimension two, the restriction
$q|_{G'}$ is an isogeny onto $E\times\bar E$ and would split the
pullback of the nonzero extension class of $G$ by that isogeny.
Injectivity of pullback on $H^1(\mathcal O)$ excludes this case.
Thus the only possible two-dimensional subgroup has
$\delta=1$, $t=0$ and trivial intersection with $\Gamma$, giving
$\operatorname{rank}\Gamma\le2$.
Finally $G'=0$ has $\delta=3$, $t=1$ and gives the same bound.
This proves it for every finitely generated $\Gamma$, hence for $S_h$.
\end{proof}

Torsion points on this graph have algebraic coordinates in both
factors but contribute rank zero. The proof supplies no third
independent algebraic graph point from algebraicity of $h$.
The nonzero algebraic-$h$ gate is consequently unchanged.
